\appendix
\section{Proof dependency graph}\label{app:depgraph}

\begin{figure}[h]
\centering
\begin{tikzpicture}[
  st/.style={draw=black!55,rounded corners=2pt,align=center,font=\scriptsize,inner sep=3pt,
             text width=2.9cm,fill=white},
  thm/.style={st,draw=black!80,fill=black!8,font=\small,text width=4.4cm},
  dep/.style={-{Stealth[length=2mm]},draw=black!70},
  ddep/.style={-{Stealth[length=2mm]},draw=black!45,dashed},
  cons/.style={st,draw=black!35,dashed,text width=9.4cm},
  lab/.style={font=\scriptsize,text=black!60}
]
\node[thm] (T) at (0,0) {\cref{thm:main}};
\node[st] (P2) at (-2.7,-1.8) {\cref{prop:reduction}};
\node[st] (P4) at (2.7,-1.8) {\cref{prop:no-concavity}};
\node[st] (L22) at (-4.5,-3.5) {\cref{lem:deficit}};
\node[st] (L23) at (-1.4,-3.5) {\cref{lem:blowup}};
\node[st] (L31) at (2.7,-3.5) {\cref{lem:pruning}};
\node[st] (L32) at (2.7,-5.2) {\cref{lem:kernel}};
\node[st] (KON) at (5.6,-5.2) {K\H{o}nig's theorem};
\node[cons] (C) at (0,-7.2)
  {\cref{cor:weighted} (weighted forms)\qquad \cref{cor:cycles} (directed cycles)\\[1pt]
   consequences of \cref{thm:main}, not used in its proof};
\draw[dep] (L22.north) -- (P2.south west);
\draw[dep] (L23.north) -- (P2.south east);
\draw[dep] (L31.north) -- (P4.south);
\draw[dep] (L32.north) -- (L31.south);
\draw[dep] (KON.north west) -- (L31.south east);
\draw[dep] (P2.north) -- (T.south west);
\draw[dep] (P4.north) -- (T.south east);
\draw[ddep] (T.south) -- (C.north);
\node[lab,anchor=west] at (0.15,-5.6) {not used above};
\end{tikzpicture}

\caption{Which numbered statement uses which. Read upwards: an arrow from $A$ to $B$ means that the
 proof of $B$ uses $A$. \emph{Look at} the two chains: the theorem uses exactly two propositions, and
 each proposition uses one substantial lemma. \emph{Why it matters:} the dashed box at the bottom shows
 that \cref{cor:weighted,cor:cycles} consume the theorem and feed nothing, so they can be skipped; and
 the two columns are independent of each other, so the reduction and the pruning can be read in
 either order.}
\label{fig:depgraph}
\end{figure}

\begin{table}[h]
\centering
\small
\renewcommand{\arraystretch}{1.2}
\begin{tabular}{@{}llp{6.6cm}@{}}
\toprule
statement & depends on & role \\
\midrule
\cref{thm:main} & \cref{prop:reduction,prop:no-concavity} & main result \\
\cref{prop:reduction} & \cref{lem:deficit,lem:blowup} & counterexample $\Rightarrow$ strict subset inequality \\
\cref{lem:deficit} & minimality of the counterexample & subset deficit in a minimal counterexample \\
\cref{lem:blowup} & \cref{lem:deficit} & amplifies the deficit by the block size \\
\cref{prop:no-concavity} & \cref{lem:pruning} & extremal argument, disjoint from \cref{prop:reduction} \\
\cref{lem:pruning} & \cref{lem:kernel}, K\H{o}nig's theorem & pruning of two families of pairs \\
\cref{lem:kernel} & --- & linear algebra of maximum-rank matrices \\
\cref{cor:weighted} & \cref{thm:main}, \cite{Seacrest2015} & weighted forms (quoted equivalence) \\
\cref{cor:cycles} & \cref{thm:main} & directed-cycle and degree consequences \\
\bottomrule
\end{tabular}
\caption{The same information as \cref{fig:depgraph}, in words.}
\label{tab:dep}
\end{table}

\section{Dictionary: notation and numbering against the original}\label{app:dictionary}

This document follows the original preprint statement by statement, but some symbols were renamed so
that no symbol carries two meanings. The table translates. The numbering of the propositions,
theorems and corollaries agrees with the original; inside the sections it does not, because remarks
are numbered in the same sequence here. The differences are: the original's Lemma~2.3 (blowup) is
this document's Lemma~\ref{lem:blowup}, the original's Lemma~3.1 (pruning) is
Lemma~\ref{lem:pruning}, and the original's Lemma~3.2 (kernels) is Lemma~\ref{lem:kernel}.

{\small
\renewcommand{\arraystretch}{1.15}
\begin{center}
\begin{tabular}{@{}llp{6.2cm}@{}}
\toprule
original & here & why the change \\
\midrule
$F,\ F_D$, $F^2$ & $F$, $F^2$ & unchanged: the vertex-set image operator \\
$F_1,\ F_2$ & $\Phi_1,\ \Phi_2$ & distinct glyphs; $F$ now has one meaning \\
$f$ (inside a block) & $\tau$ & avoids the $F/f$ collision \\
$T\subseteq E$ & $Y\subseteq E$ & $T$ is reserved for the tournament $T_m$ \\
$D_T$ & $D_Y$ & same reason \\
$Q=F(\{v\})\cap(E\setminus T)$ & $Q_v$ & subscripted: $Q$ alone is a family of pairs in \cref{sec:incompatibility} \\
$C=F(Q)\setminus(I\cup T\cup Q\cup F(T))$ & $\Lambda$ & $C$ alone is the right family in \cref{lem:pruning} \\
$d_1,\ d_2$ & $d_1,\ d_2$ & unchanged (degrees at a vertex) \\
$d$ (uncovered), $d'$ & $u$, $u'$ & $d$ is now only a degree \\
$g(T)$ & $g(Y)$ & notation only \\
$H_0=F(S)\cup F^2(S)$ & $W$ & $H$ is the collision set in \cref{lem:pruning} \\
$R,\ C$ (families) & $\mathcal L,\ \mathcal R$ & ``left'' and ``right'' is what they are \\
$R,\ C$ (complements in \S4) & $\mathcal L,\ \mathcal R$ & the same role after pruning \\
$P',\ Q'$ & $\mathcal L',\ \mathcal R'$ & primes never attach to the extremal pair \\
$Z,\ H$ & $Z,\ H$ & unchanged: conflict points, collisions \\
$Z(v)$ (non-neighbors, \S6) & $\overline N(v)$ & $Z$ is taken by the conflict points \\
$u^r,\ v^c$ (lifts) & $\lambda^r,\ \kappa^c$ & $u$ is the uncovered count in \S4 \\
$\lambda^r_a,\ \mu^c_b$ (coefficients) & $\xi^r_a,\ \eta^c_b$ & frees $\lambda$ for the lift vector \\
$k,\ \ell,\ t$ & $k,\ \ell,\ t$ & unchanged \\
$M$ (maximum matching) & $\mathcal M$ & $M=\abs P+\abs Q$ in \S4 \\
\bottomrule
\end{tabular}
\end{center}}

\section{Verification notes}\label{app:checks}

This section records what was checked and how, so that a reader can decide what to trust and what to
re-check. Three kinds of check were used.

\paragraph{(V1) Re-derivation.} Every display of the original was re-derived from the definitions
during the rewrite, and every inequality direction was checked at the point of use. Three places
deserve a note because they are easy to misread.
\begin{itemize}[leftmargin=1.4em,itemsep=2pt,topsep=2pt]
\item In the ledger step of \cref{lem:deficit} the second out-neighborhood of $v$ is confined to $I\cup Y\cup F(Y)$
 \emph{after} the deletion, and this uses $v\in S$ (which comes from the choice of $v$, not from the
 definition of $Q_v$). Without $v\in S$ the inclusion fails.
\item In the block accounting in \cref{lem:blowup} the sum over $S\setminus W$ is extended to $S\setminus F(S)$; this is an
 inequality, not an identity, because a block in $F(S)$ is already fully counted. Treating it as an
 identity would be a mistake.
\item In the kernel step of \cref{lem:pruning} the auxiliary matrix $D^{p,i}$ has maximum rank on $T^{p,i}$, and the
 hypothesis $(s,j)\in T^{p,i}$ is exactly $s\to j$; the vanishing is asserted only for such pairs.
 For $(s,j)\notin T^{p,i}$ the entry $b_{sj}$ is zero and there is nothing to prove.
\end{itemize}

\paragraph{(V2) Exhaustive and random computation.} The combinatorial statements of the paper were
checked on small instances by the scripts in \code{checks/verify.py} of the repository. All checks
pass. The checks are implications and identities rather than the statements themselves, since the
hypotheses of \cref{lem:deficit,lem:blowup,prop:no-concavity} are (conjecturally) unsatisfiable ---
that is the theorem.
\begin{center}
\small
\renewcommand{\arraystretch}{1.15}
\begin{tabular}{@{}p{8.5cm}p{5.6cm}@{}}
\toprule
what & result \\
\midrule
\cref{thm:main} on all oriented graphs with $n\le4$ (760 graphs) and $20000$ random graphs with $5\le n\le9$
 & holds in every case \\
the ledger of the ledger step of \cref{lem:deficit}: the four inclusions, \eqref{eq:g-diff}, the sandwich, and the strict increment
 & $2336$ random $(D,S,Y)$ configurations, no failure \\
the algebra of \cref{lem:blowup}: \eqref{eq:FG}, \eqref{eq:FG2}, \eqref{eq:blowup-main}
 & $24942$ random $(D,m,U)$ instances, no failure \\
the identities of \cref{prop:no-concavity}: fibrewise images, complement gain, diagonal survival, $H$ uncovered
 & $250$ random compatible pairs, no failure \\
the statement of \cref{lem:pruning}, by brute force over all subfamilies
 & $1200$ random instances with $\abs X\le4$, no counterexample \\
\cref{cor:cycles}(i) and the partition \eqref{eq:partition} on $3$-cycle-free graphs; (ii) on random graphs; (iii) on regular girth-$4$ graphs
 & $34114$ graphs, no failure \\
\bottomrule
\end{tabular}
\end{center}

\paragraph{(V3) Formalisation.} A Lean~4 development accompanies this document (\code{lean/} in the
repository). It contains the definitions of oriented graph, first and second out-neighborhood, the
statements of \cref{thm:main,prop:reduction,prop:no-concavity,lem:pruning}, a fully checked proof of
the final assembly of \cref{thm:main} from the two propositions, and a fully checked proof of
\cref{cor:cycles} from \cref{thm:main}. The deep lemmas (\cref{lem:deficit,lem:blowup,lem:pruning})
are stated in Lean but their proofs are not formalised; they are recorded as explicit hypotheses so
that the reader can see exactly which part of the argument is machine-checked and which is not. The
formalisation therefore checks the \emph{statements} and the \emph{logical architecture}, not the
three engines.

\paragraph{What this does not establish.} A pleasant presentation and a passing test suite do not
prove a theorem. The only complete proof in this document is the mathematical text itself; the
computations of (V2) test small cases, and the formalisation of (V3) checks the statements and the
architecture. Neither substitutes for reading the proofs of \cref{lem:deficit,lem:blowup,lem:pruning}.

\section{Index of statements and figures}\label{app:index}

{\small
\renewcommand{\arraystretch}{1.12}
\begin{center}
\begin{tabular}{@{}llp{7.6cm}@{}}
\toprule
\cref{thm:main} & main theorem & every nonempty finite oriented graph has $v$ with $\abs{N_1^+(v)}\le\abs{N_2^+(v)}$ \\
\cref{prop:reduction} & reduction & a counterexample yields \eqref{eq:strict-growth} with positive indegrees \\
\cref{lem:deficit} & subset deficit & $\abs{F^2(S)\setminus F(S)}<\abs{F(S)\setminus S}$ in a minimal counterexample \\
\cref{lem:blowup} & blowup & transitive-tournament blowup gives \eqref{eq:strict-growth} \\
\cref{def:conflict} & conflicts & conflict, conflict point, collision, covering, compatible \\
\cref{lem:pruning} & pruning & compatible subfamilies losing at most $\abs H+q$ members \\
\cref{lem:kernel} & rank lemma & kernels at maximum support rank have disjoint supports \\
\cref{prop:no-concavity} & no graph & \eqref{eq:impossible} is impossible \\
\cref{cor:weighted} & weighted forms & vertex- and arc-weighted statements \\
\cref{cor:cycles} & directed cycles & degree and order bounds without directed triangles \\
\bottomrule
\end{tabular}
\end{center}}

\begin{center}
\small
\renewcommand{\arraystretch}{1.1}
\begin{tabular}{@{}ll@{}}
\cref{fig:definitions}: the two neighborhoods & \cref{fig:roadmap}: architecture of the proof \\
\cref{fig:deficit}: the split of $F(S)$ & \cref{fig:deletion}: the deletion step \\
\cref{fig:blowup}: the blowup & \cref{fig:conflict}: a conflict as a rectangle \\
\cref{fig:pruning-graph}: the augmented graph & \cref{fig:cover}: the charge to the cover \\
\cref{fig:fourmaps}: the commutative square & \cref{fig:grid}: families in the grid \\
\cref{fig:diagonal}: the diagonal survives & \cref{fig:extremal}: the four moves of \cref{sec:incompatibility} \\
\cref{fig:beta}: the arc weighting & \cref{fig:cycles}: non-neighbors and the partition \\
\cref{fig:depgraph}: dependencies & \\
\end{tabular}
\end{center}
